\begin{center}\begin{tikzpicture}[font=\small,>=stealth,line join=round]
\draw[gray!40,rounded corners=3pt] (0,0) rectangle (14.2,9.0);
\node at (7.1,8.4) {例6：去掉奇的部分，左右才会等高};
\begin{scope}[shift={(3.6,3.3)},x=1.05cm,y=1.65cm]\draw[->,gray] (-2.5,0)--(2.7,0) node[right]{$x$};\draw[->,gray] (0,-0.2)--(0,2.5) node[left]{$y$};\begin{scope}\clip (-2.5,-0.2) rectangle (2.5,2.3);\draw[teal!80!black,thick,domain=-2.5:2.5,samples=120] plot(\x,{1});\draw[blue,thick,domain=-2.5:2.5,samples=120] plot(\x,{1+\x/(\x*\x+1)});\end{scope}\fill[blue] (-1,.5) circle(2pt) node[below left]{$(-1,1/2)$};\fill[blue] (1,1.5) circle(2pt) node[above right]{$(1,3/2)$};\end{scope}\node[align=left,anchor=west] at (7.8,4.8) {$f(x)=1+\dfrac{ax}{x^2+1}$\\[14pt]\textcolor{blue}{$a=1$：两边值不同}\\[12pt]\textcolor{teal!80!black}{$a=0$：恒为 $1$，是偶函数}\\[12pt]要对所有 $x$ 都相同，只能 $a=0$};
\node[text width=12.8cm,align=center] at (7.1,.65) {绿色恒等于1；蓝色为参数 $a=1$ 的反例。纵轴显示实际函数值。};
\end{tikzpicture}\end{center}
